\documentclass[11pt]{beamer}

% ============================================================
% Theme & Appearance
% ============================================================
\usetheme{Madrid}
\usecolortheme{seahorse}
\setbeamertemplate{navigation symbols}{}

% ============================================================
% Encoding & Fonts (pdfLaTeX-safe)
% ============================================================
\usepackage[utf8]{inputenc}
\usepackage[T1]{fontenc}
\usepackage{lmodern}

% ============================================================
% Math & Graphics
% ============================================================
\usepackage{amsmath, amssymb}
\usepackage{graphicx}
\usepackage{tikz}
\usetikzlibrary{arrows.meta}

% ============================================================
% Code Listings
% ============================================================
\usepackage{listings}
\usepackage{xcolor}

\lstset{
	basicstyle=\ttfamily\small,
	frame=single,
	breaklines=true,
	columns=fullflexible,
	keywordstyle=\color{blue},
	commentstyle=\color{gray},
	stringstyle=\color{red},
	showstringspaces=false
}

% ============================================================
% Title Information
% ============================================================
\title{Concurrency \& Distributed Systems}
\subtitle{Week 1 --- Why Concurrency Exists (Pain Demo Part 1)}
\author{Dr. Mohamad Aoude}
\institute{Lebanese University \\ Faculty of Engineering}
\date{\today}

% ============================================================
\begin{document}
	% ============================================================
	
	\begin{frame}
		\titlepage
	\end{frame}
	
	\section{The Problem}
	
	% ------------------------------------------------------------
	\begin{frame}[fragile]{The Architecture We Start With}
	\textbf{Single-threaded server}
		
	\begin{lstlisting}[language=Java]
		while (true) {
			Socket client = serverSocket.accept();
			handleRequest(client);
		}
	\end{lstlisting}

		
		\textbf{Properties:}
		\begin{itemize}
			\item Processes one request at a time
			\item New clients must wait
		\end{itemize}
		

		
		\begin{block}{Question}
			Is this architecture wrong?
		\end{block}
	\end{frame}
	
	% ------------------------------------------------------------
	\begin{frame}[fragile]{Inside handleRequest()}
	\begin{lstlisting}[language=Java]
		private static void handleRequest(Socket client) {
			Thread.sleep(100); // simulate DB call
			sendResponse(client);
		}
	\end{lstlisting}
		

		
		\textbf{Important:}
		\begin{itemize}
			\item 100ms is simulated I/O wait
			\item CPU is idle during sleep
		\end{itemize}
		

		
		\begin{alertblock}{Key Insight}
			The thread is blocked. The CPU is mostly unused.
		\end{alertblock}
	\end{frame}
	
	% ------------------------------------------------------------
	\begin{frame}{Experiment 1: One Client}
		\textbf{Measured:}
		\begin{itemize}
			\item Response time $\approx 100$ms
			\item CPU usage $\approx 0$--$2\%$
		\end{itemize}
		

		
		\[
		\text{CPU Utilization} \approx \frac{2ms}{100ms} = 2\%
		\]
		

		
		\begin{alertblock}{Interpretation}
			Server is idle about 98\% of the time.
		\end{alertblock}
	\end{frame}
	
	% ------------------------------------------------------------
	\begin{frame}{Experiment 2: 5 Concurrent Clients}
		All arrive at $t=0$.
		

		
		Sequential processing:
		\[
		5 \times 100ms = 500ms
		\]
		

		
		\textbf{Observed:}
		\begin{itemize}
			\item Latency $\approx 500$ms
			\item CPU still $\approx 2\%$
		\end{itemize}
		

		
		\begin{block}{Critical Question}
			Why does latency multiply if CPU is not busy?
		\end{block}
	\end{frame}
	
	% ------------------------------------------------------------
	\begin{frame}{Experiment 3: 50 Concurrent Clients}
		Sequential prediction:
		\[
		50 \times 100ms = 5000ms
		\]
		

		
		\textbf{Observed:}
		\begin{itemize}
			\item Latency $\approx 5$ seconds
			\item CPU mostly idle
		\end{itemize}
		

		\begin{alertblock}{This is I/O-bound death}
			Hardware is not the problem. Architecture is the problem.
		\end{alertblock}
	\end{frame}
	
	% ------------------------------------------------------------
	\begin{frame}{Visualizing the Queue}
		\begin{center}
			\begin{tikzpicture}[x=0.15cm,y=0.8cm,>=Stealth]
				\draw[->] (0,0) -- (65,0) node[right]{time};
				
				\node[anchor=east] at (0,1) {Client 1:};
				\node[anchor=east] at (0,2) {Client 2:};
				\node[anchor=east] at (0,3) {Client 3:};
				\node[anchor=east] at (0,5) {Client 50:};
				
				\draw[thick] (2,1) rectangle +(6,0.4);
				\draw[thick] (7,2) rectangle +(6,0.4);
				\draw[thick] (12,3) rectangle +(6,0.4);
				\node at (15,4) {$\vdots$};
				\draw[thick] (52,5) rectangle +(6,0.4);
			\end{tikzpicture}
		\end{center}
		

		
		All requests arrive simultaneously. Only one is processed at a time.
	\end{frame}
	
	% ------------------------------------------------------------
	\begin{frame}{Mathematical Framing}
		Each request:
		\[
		T_{compute} \approx 2ms,\quad
		T_{wait} \approx 98ms,\quad
		T_{service} = 100ms
		\]
		

		
		CPU utilization:
		\[
		\frac{T_{compute}}{T_{service}} \approx \frac{2}{100} = 2\%
		\]
		

		\begin{block}{Conclusion}
			The system is not CPU-bound. It is waiting-bound.
		\end{block}
	\end{frame}
	
	% ------------------------------------------------------------
	\begin{frame}{The Fundamental Question}
		\Large
		If the CPU is idle about 98\% of the time,\par

		why do we allow it to serve only one request?\par
		

		\vspace{4mm}
		\begin{alertblock}{This question creates Concurrency.}
		\end{alertblock}
	\end{frame}
	% ------------------------------------------------------------
	\begin{frame}[fragile]{Student Instructions: How to Run (Two Terminals)}
		\textbf{Open two terminals.}
		
		\vspace{2mm}
		\textbf{Terminal 1 --- Start Server}
		\begin{lstlisting}
			javac SingleThreadedServer.java
			java SingleThreadedServer
		\end{lstlisting}
		

		\textbf{Terminal 2 --- Run Client}
		\begin{lstlisting}
			javac LoadClient.java
			java LoadClient
		\end{lstlisting}
		

		\begin{block}{Rule}
			Keep the server running. Only rerun the client between tests.
		\end{block}
	\end{frame}
	
	% ------------------------------------------------------------
	\begin{frame}[fragile]{Student Instructions: Run 3 Experiments}
		In \texttt{LoadClient.java}, change:
		
		\begin{lstlisting}[language=Java]
			int concurrentClients = 1;   // Test A
		\end{lstlisting}
		

		Then rerun:
		\begin{lstlisting}
			javac LoadClient.java
			java LoadClient
		\end{lstlisting}
		

		Repeat for:
		\begin{lstlisting}[language=Java]
			int concurrentClients = 5;   // Test B
			int concurrentClients = 50;  // Test C
		\end{lstlisting}
		

		\begin{alertblock}{Important}
			Do NOT change the server. Only change the client concurrency.
		\end{alertblock}
	\end{frame}
	
	% ------------------------------------------------------------
	\begin{frame}{Student Instructions: What to Record}
		For each run (1, 5, 50 clients), record:
		
		\begin{itemize}
			\item \textbf{Average latency (ms)} (or the value printed by your client)
			\item \textbf{Worst latency (ms)} if available
			\item \textbf{CPU usage} (Task Manager / \texttt{htop})
			\item \textbf{One sentence:} what happened and why?
		\end{itemize}
		

		\begin{block}{Deliverable}
			A small table with 3 rows (1, 5, 50) + your explanation.
		\end{block}
	\end{frame}
	
	% ------------------------------------------------------------
	\begin{frame}{Expected Trend (So You Know You're Correct)}
		\begin{itemize}
			\item 1 client: latency $\approx 100$ms
			\item 5 clients: latency $\approx 500$ms
			\item 50 clients: latency $\approx 5000$ms
		\end{itemize}
		

		\begin{alertblock}{Key Observation}
			CPU stays low: the system is waiting-bound, not CPU-bound.
		\end{alertblock}
	\end{frame}
	
\end{document}
